\par\medskip\begin{small}\noindent\textbf{Quellenhinweis.} Die Anzeige vom 20.~April 1816 ist nach \textit{Werke}, Band~IV (1873), S.~364--368, in modernes Deutsch übertragen. Das Titelblatt und S.~24 von \href{https://gdz.sub.uni-goettingen.de/id/PPN603351808}{Johann Christoph Schwabs \textit{Commentatio} (Stuttgart 1814)} wurden im Exemplar der SUB Göttingen, 8 MATH III, 65, eingesehen. Die lateinischen Zitate folgen der Anzeige von Gauss. Matthias Metternichs Werk wird nach der Mainzer Ausgabe von 1815 zitiert. Francesco Valentinis \href{https://doi.org/10.3931/e-rara-91963}{deutsch-italienisches Wörterbuch} belegt die mathematische Bedeutung von \textit{angebliche Größe}, hier als \textit{angebbare Größe} wiedergegeben.\par\end{small}
